% Nová příprava: 2026-10-28; 1. ročník SPŠ, etapa 13–24.
% Rozpis: 1_rocnik/tematicke_plany/Rozpis_vyuky_MAT_1_rocnik.tex, hodina 13.
% Původ: vlastní formulace výkladu a vlastní úlohy; metodická inspirace uvedenými zdroji.
% Kontrola, obsahové porovnání zdrojů a časy: planovani/DOKONCENI_PRIPRAV_13_24.md.
% Učebnice: ucebnice/01_Zakladni_poznatky/Mnoziny/Množiny.tex
% oddíly Rovnost množin, Podmnožiny, Množinové operace a Vennovy diagramy; dle tématu hodiny.
% Realisticky.cz: externi_zdroje/realisticky/01_Zakladni_poznatky/03_Mnoziny/010301_Mnoziny_I/010301_Lekce.pdf
% Realisticky.cz: externi_zdroje/realisticky/01_Zakladni_poznatky/03_Mnoziny/010302_Mnoziny_II_P/010302_Lekce.pdf
\documentclass[11pt,a4paper]{article}
% Spolecny styl pro ucitelske pripravy MAT (Work/MAT)
% Verze 1.1, 2026-10-23
\NeedsTeXFormat{LaTeX2e}
\ProvidesPackage{priprava}[2026/10/23 v1.1 Sablona ucitelskych priprav MAT]

% Zakladni typografie a cestina
\RequirePackage[T1]{fontenc}
\RequirePackage[utf8]{inputenc}
\RequirePackage[czech]{babel}
\RequirePackage{lmodern}
\RequirePackage{microtype}
\RequirePackage{amsmath,amssymb}
\RequirePackage{xcolor}
\RequirePackage{enumitem}
\RequirePackage[most]{tcolorbox}
\RequirePackage{tabularx,booktabs,array}
\RequirePackage{fancyhdr}
\RequirePackage{needspace}
\RequirePackage{geometry}
\RequirePackage[hidelinks]{hyperref}

\geometry{a4paper,top=17mm,bottom=18mm,left=23mm,right=23mm}
\definecolor{notegray}{gray}{0.39}
\setlength{\parindent}{0pt}
\setlength{\parskip}{2.5pt plus 1pt minus 0.5pt}
\setlist[itemize]{topsep=2pt,itemsep=1pt,parsep=0pt,leftmargin=1.45em}
\setlist[enumerate]{topsep=2pt,itemsep=2pt,parsep=0pt,leftmargin=1.65em}
\setcounter{secnumdepth}{0}
\renewcommand{\arraystretch}{1.12}

% Zahlavi: vlevo rocnik | globalni tema, vpravo cislo: nazev hodiny
% Zapati zustava zcela prazdne.
\pagestyle{fancy}
\fancyhf{}
\renewcommand{\headrulewidth}{0.35pt}
\renewcommand{\footrulewidth}{0pt}
\setlength{\headheight}{14pt}
\newcommand{\nastavPripravu}[4]{%
  \fancyhead[L]{\small #1\enspace\textbar\enspace #2}%
  \fancyhead[R]{\small #3: #4}%
}

% Minimalisticke cernobile boxy. Kazda definice zvlast.
\tcbset{commonlesson/.style={enhanced,breakable=false,colback=white,colframe=black,
  colbacktitle=white,coltitle=black,boxrule=0.45pt,
  arc=0.8mm,outer arc=0.8mm,left=3mm,right=3mm,top=1.6mm,bottom=1.6mm,
  toptitle=1mm,bottomtitle=0.7mm,fonttitle=\bfseries,
  attach boxed title to top left={xshift=2.5mm,yshift=-1.15mm},
  boxed title style={colback=white,colframe=white,boxrule=0pt,left=1mm,right=1mm}}}
\newtcolorbox{zadaniBox}[1][]{commonlesson,title={Zadání},#1}
\newtcolorbox{prikladBox}[1][]{commonlesson,title={Příklad},#1}
\newtcolorbox{definiceBox}[1][]{commonlesson,title={Definice},#1}

% Spolecny box pro zadani a jeho reseni, bez mezery mezi castmi.
% Pouziti:
%   \begin{prikladReseniBox}
%     Zadani ulohy.
%     \tcblower
%     \begin{align*} ... \end{align*}
%   \end{prikladReseniBox}
% Spodni cast ma jemny sedy podklad a delici linku; lze v ni sazet
% vice radku matematiky. Obe casti maji jen jeden spolecny ramecek.
\tcbset{solutionlesson/.style={
  commonlesson,
  enhanced jigsaw,
  breakable,
  colbacklower=black!3,
  segmentation style={solid,draw=black!25,line width=0.35pt},
  fontlower=\small\color{notegray},
  before lower={\noindent\textbf{Řešení:}\enspace},
  leftlower=3mm,rightlower=3mm
}}
\newtcolorbox{prikladReseniBox}[1][]{solutionlesson,title={Příklad},#1}
\newtcolorbox{zadaniReseniBox}[1][]{solutionlesson,title={Zadání},#1}

% Pomucky pro autora pripravy: sedive poznamky; \reseni pro starsi pripravy
\newcommand{\poznamka}[1]{\par\begingroup\small\color{notegray}\itshape
  \textnormal{\textbf{Poznámka pro učitele.}}\enspace #1\par\endgroup}
\newcommand{\reseni}[1]{\par\begingroup\small\color{notegray}
  \textbf{Očekávané řešení:}\enspace #1\par\endgroup}
\newcommand{\blok}[1]{\needspace{5\baselineskip}\subsection*{#1}}
\newcommand{\vyrok}[1]{\uv{#1}}

\nastavPripravu{1. ročník}{Množiny}{13}{Podmnožina a rovnost}
\begin{document}
\blok{Od jednotlivého prvku k celé množině}
V minulé hodině jsme rozhodovali, zda určitý objekt patří do množiny.
Dnes porovnáváme celé množiny: leží všechny prvky jedné také ve druhé?
Uvažuj $A=\{2,5\}$ a $B=\{0,2,5,8\}$. Nejdříve ověř každý prvek $A$ zvlášť.

\begin{definiceBox}
Množina $A$ je \textbf{podmnožinou} množiny $B$, právě když každý prvek $A$
patří také do $B$. Píšeme $A\subseteq B$, tedy
$\forall x:\ x\in A\Rightarrow x\in B$.
\end{definiceBox}
Zápis $A\nsubseteq B$ znamená, že existuje alespoň jeden prvek $A$,
který do $B$ nepatří. K vyvrácení stačí tento konkrétní prvek; k potvrzení
konečného výčtu musíme ověřit všechny prvky.

\begin{prikladReseniBox}[title={P1: oběma směry},breakable=false]
Pro uvedené $A,B$ rozhodni o $A\subseteq B$ a $B\subseteq A$.
\tcblower
$A\subseteq B$: prvky $2$ i $5$ jsou v $B$. Obrácený vztah neplatí,
například $0\in B$, ale $0\notin A$. Podmnožina tedy není souměrný vztah.
\end{prikladReseniBox}

\begin{definiceBox}
Množiny $A,B$ jsou \textbf{rovny}, mají-li právě stejné prvky.
Platí $A=B$ právě tehdy, když $A\subseteq B$ a současně $B\subseteq A$.
\end{definiceBox}
\begin{definiceBox}
$A$ je \textbf{vlastní podmnožinou} $B$, jestliže $A\subseteq B$ a $A\ne B$.
Značíme $A\subsetneq B$.
\end{definiceBox}
Rovnost se nemění přeuspořádáním výčtu ani opakováním prvku.
Značku $\subset$ různí autoři používají různě; zde pro jednoznačnost rozlišujeme
$\subseteq$ a $\subsetneq$.

\begin{prikladReseniBox}[title={P2: stejná čísla, různé zápisy},breakable=false]
Porovnej $C=\{x\in\mathbb Z\mid -2<x\le2\}$,
$D=\{2,0,-1,1,2\}$ a $E=(-2;2\rangle$.
\tcblower
$C=\{-1,0,1,2\}=D$. Oba zápisy dávají právě stejné čtyři prvky.
Platí $C\subsetneq E$: všechna čtyři čísla leží v intervalu,
ale například $\frac12\in E$ a $\frac12\notin C$. Rozhoduje i číselný obor.
\end{prikladReseniBox}

\blok{Prázdná množina a význam závorek}
Vždy $A\subseteq A$. Také $\varnothing\subseteq A$: neexistuje prvek prázdné
množiny, který by vztah porušil. Netvrdíme tím $\varnothing\in A$.
O prvku píšeme $\in$, o zahrnutí celé množiny $\subseteq$.

\begin{prikladReseniBox}[title={P3: co je objekt a co soubor?},breakable=false]
Pro $M=\{0,2\}$ rozhodni: $0\in M$, $\{0\}\subseteq M$,
$\{0\}\in M$, $\varnothing\subseteq M$ a $\varnothing\in M$.
\tcblower
Postupně pravda, pravda, nepravda, pravda, nepravda.
Prvky $M$ jsou čísla $0$ a $2$, nikoli množiny $\{0\}$ nebo $\varnothing$.
Množina $\{0\}$ obsahuje jen číslo $0$, které v $M$ skutečně je.
\end{prikladReseniBox}
\poznamka{Nejdříve nech vyslovit definici vlastními slovy. Jádro tvoří P1–P3
a S1; rozlišení vnořených množin dále nerozvíjej. Počítej se třemi minutami
organizace a přibližně 39 minutami práce a kontroly, závěr a domácí zadání 3 minuty.}

\blok{Samostatné procvičení}
\begin{zadaniReseniBox}[title={S1: rozhodnutí musí mít důvod}]
\begin{enumerate}[label=\alph*)]
\item Pro $A=\{-1,0,4\}$, $B=\{0,4\}$, $C=\{4,0\}$ zapiš všechny platné
vztahy $\subseteq$ mezi různě označenými množinami. Které jsou si rovny?
\item Platí $\{0\}=\varnothing$? Uveď rozlišující prvek.
\item Jsou si rovny $\{x\in\mathbb N\mid x<4\}$ a $\{x\in\mathbb Z\mid0<x<4\}$?
\item Platí $\mathbb N\subseteq\mathbb Z$ a $\mathbb Z\subseteq\mathbb N$?
Používáme $\mathbb N=\{1,2,\ldots\}$.
\item Vypiš všechny podmnožiny $\{p,q\}$, kde $p\ne q$. Které jsou vlastní?
\end{enumerate}
\tcblower
\begin{enumerate}[label=\alph*)]
\item $B\subseteq A$, $C\subseteq A$, $B\subseteq C$, $C\subseteq B$.
$B=C$; $A$ se jim nerovná kvůli prvku $-1$.
\item Ne: $0\in\{0\}$, ale $0\notin\varnothing$.
\item Ano, obě jsou $\{1,2,3\}$. Nula první množině nepatří.
\item První vztah ano, druhý ne, například kvůli $0\in\mathbb Z$.
\item $\varnothing$, $\{p\}$, $\{q\}$, $\{p,q\}$.
Vlastní jsou první tři, poslední je podmnožinou sebe samé.
\end{enumerate}
\end{zadaniReseniBox}

\blok{Volitelná rezerva}
\begin{zadaniReseniBox}[title={R1: systematický výčet a obecný vztah}]
\begin{enumerate}[label=\alph*)]
\item Vypiš všechny podmnožiny $\{1,4,6\}$ podle počtu prvků.
\item Zdůvodni: z $A\subseteq B$ a $B\subseteq C$ plyne $A\subseteq C$.
\end{enumerate}
\tcblower
a) $\varnothing$; $\{1\},\{4\},\{6\}$;
$\{1,4\},\{1,6\},\{4,6\}$; $\{1,4,6\}$. Je jich osm.
Každý ze tří prvků buď zařadíme, nebo nezařadíme.

b) Libovolný prvek $A$ je podle první podmínky v $B$ a podle druhé také v $C$.
Tím jsme ověřili definici $A\subseteq C$, nikoli jen jeden číselný pokus.
\end{zadaniReseniBox}
\poznamka{Rezerva přibližně 8–12 minut, nepočítá se navíc do jádra.
Závěrem požaduj protipříklad k jednomu neplatnému zahrnutí.}

\blok{Domácí procvičení}
\begin{zadaniReseniBox}[title={D1: přibližně 12–15 minut}]
\begin{enumerate}[label=\alph*)]
\item Pro $K=\{-2,0,3\}$, $L=\{3,-2,0,0\}$ a $M=\{-2,3\}$ urči rovnost a vlastní zahrnutí.
\item Rozhodni a zdůvodni: $\varnothing\subseteq\{0\}$, $0\in\varnothing$,
$\{3\}\subseteq K$.
\item Zapiš výčtem $P=\{x\in\mathbb Z\mid-1\le x<3\}$ a ověř $P\subseteq\mathbb N$.
\item Vypiš všechny podmnožiny $\{-1,2\}$.
\item Najdi dvě různé dvouprvkové množiny, z nichž žádná není podmnožinou druhé.
\end{enumerate}
\tcblower
a) $K=L$, $M\subsetneq K$ a $M\subsetneq L$.
b) Pravda (žádný odporující prvek), nepravda (prázdná množina nemá prvky),
pravda ($3\in K$).
c) $P=\{-1,0,1,2\}$; zahrnutí neplatí například kvůli $-1$.
d) $\varnothing,\{-1\},\{2\},\{-1,2\}$.
e) Například $\{1,2\}$ a $\{2,3\}$: rozhodují prvky $1$ a $3$.
\end{zadaniReseniBox}
\end{document}
